\section{从特征工程到端到端学习}

\subsection{传统特征表示}

在深度学习时代之前（约 2000--2012 年），计算机视觉的标准范式是\textbf{手工特征提取 + 线性分类器}。研究者会设计各种特征表示来捕获图像中的关键信息：

\begin{figure}[H]
\centering
\includegraphics[width=0.95\textwidth]{slides-images/slide-012.jpg}
\caption{传统特征表示：颜色直方图（左上）和方向梯度直方图 HoG（右侧）\protect\footnotemark}
\end{figure}
\footnotetext{Slides 第13页。}

\begin{itemize}
    \item \textbf{颜色直方图}（Color Histogram）：将色彩空间离散化为若干桶，统计每个桶中像素数量。捕获颜色分布，但完全丢弃空间结构。
    \item \textbf{方向梯度直方图}（HoG, Histogram of Oriented Gradients）：在图像的每个局部区域计算边缘方向的分布。捕获空间结构，但丢弃颜色信息。
    \item \textbf{拼接特征}：实践中通常将多种特征向量拼接起来，形成一个大的特征表示。
\end{itemize}

\subsection{特征提取 vs 端到端学习}

\begin{figure}[H]
\centering
\includegraphics[width=0.95\textwidth]{slides-images/slide-015.jpg}
\caption{两种范式对比：（A）手工特征 + 学习分类器 vs（B）端到端神经网络\protect\footnotemark}
\end{figure}
\footnotetext{Slides 第16页。}

\begin{importantbox}{端到端学习的核心理念}
\textbf{梯度下降可能比你更擅长编程，海量数据可能比你更了解问题。}

传统方法中，特征提取部分由人类设计（可能包含错误直觉或遗漏），只有分类器部分从数据中学习。而端到端神经网络让\textbf{整个系统}——从原始像素到最终预测——都通过梯度下降从数据中学习。

过去十余年的实践反复证明：端到端学习在几乎所有问题上都优于手工特征工程。
\end{importantbox}

\begin{knowledgebox}{人类设计者的角色变化}
在端到端学习时代，人类的角色从``设计特征提取器''变为``设计网络架构''。两者的区别是：
\begin{itemize}
    \item 设计特征提取器 = 定义一个\textbf{具体函数}
    \item 设计网络架构 = 定义一个\textbf{函数族}（由计算图结构决定），具体参数留给梯度下降从数据中学习
\end{itemize}
\end{knowledgebox}

\subsection{本章小结}

从手工特征提取到端到端学习是计算机视觉发展的关键转折。端到端学习的成功建立在三个基础之上：（1）足够强大的函数族（深度神经网络），（2）足够多的数据（ImageNet 等大规模数据集），（3）足够快的计算（GPU）。接下来我们将学习专门为图像设计的计算图元件——卷积层和池化层。

%% ========== Part 3: 卷积神经网络 ==========

